\documentclass[11pt]{article}
\usepackage[margin=1in]{geometry}
\usepackage{amsmath,amssymb,amsfonts}
\usepackage{booktabs,longtable,array}
\usepackage{xcolor}
\usepackage{hyperref}
\usepackage{enumitem}
\hypersetup{colorlinks=true,linkcolor=blue,urlcolor=blue,citecolor=blue}

\title{The W33/Q4 Diamond Machine:\\A Claim-Stratified Master Synthesis}
\author{Wil Dahn\\W33 Theory Project}
\date{June 19, 2026}

\newcommand{\Qfour}{Q_4}
\newcommand{\W}{\mathcal W(3,3)}
\newcommand{\Fthree}{\mathbb F_3}
\newcommand{\claim}[2]{\noindent\textbf{[#1]}\;#2}

\begin{document}
\maketitle

\begin{abstract}
This document merges the audited Q4 diamond paper with the broader Q4 Diamond Machine synthesis. The guiding rule is claim stratification: exact arithmetic, executable certificates, structural proof notes, simulation gates, engineering gates, optimization frontiers, architecture boundaries, and speculative bridges are separated rather than collapsed into one theorem. The merged picture keeps the machine-audited Q4 corrections - especially the rolling-epoch repair and the new [[32,4,4]] gauge-quotient certificate - while importing the W(3,3), heptad/Reye, Kagome, packet-runtime, non-Clifford port, and chip-falsifier layers as explicitly labeled claims.
\end{abstract}

\tableofcontents
\newpage

\section{Claim Classes}

\begin{center}
\begin{tabular}{lll}
\toprule
Label & Meaning & Required evidence\\
\midrule
EXACT & finite arithmetic or finite graph computation & script/result certificate\\
CERT & executable code certificate & tool plus JSON output\\
STRUCT & proof-note theorem or construction & proof note, assumptions named\\
SIM & synthetic or numerical simulator & model, seed, pass/fail tolerances\\
ENG & engineering feasibility gate & physical assumptions and margins\\
OPT & optimization frontier & exact objective plus witness/bound status\\
ABI & architecture boundary/interface & packet contract and verification signature\\
SPEC & speculative bridge & not theorem-grade until audited\\
\bottomrule
\end{tabular}
\end{center}

\section{Audited Q4 Spine}

\claim{EXACT}{The 4-cube has \(|V(\Qfour)|=16\), \(|E(\Qfour)|=32\), and 24 square faces.}

\claim{CERT}{BT1327 audits the Q4 number table and BT1328 repairs the master epoch. The corrected epoch is}
\[
3660=6\cdot540+180,\qquad 3\cdot180=540,\qquad 3\cdot3660=10980.
\]
Thus \(10980\) is a three-frame rolling chart-phase closure, not \(\mathrm{lcm}(3660,1620)\), whose literal value is \(98820\).

\claim{CERT}{BT1338 extracts the literal cubical Q4 chain checks and finds \(k=0\) for the raw contractible cube:}
\[
\operatorname{rank}(\partial_1)=15,\qquad \operatorname{rank}(\partial_2)=17,
\qquad 32-15-17=0.
\]
This prevents an overclaim: the raw cube is not itself the \([[32,4,4]]\) code.

\claim{CERT}{BT1341 supplies the missing gauge quotient. Four global flux functionals on the 17-dimensional cycle space have a 13-dimensional kernel used as the Z-check space. With the 15 vertex/star X-checks,}
\[
 n=32,
 \qquad \operatorname{rank}(H_X)=15,
 \qquad \operatorname{rank}(H_Z)=13,
 \qquad k=4.
\]
The certificate also finds \(d_X=d_Z=4\), so the quotient object is \([[32,4,4]]\).

\claim{CERT}{BT1344 canonicalizes this quotient under the 384 cube automorphisms. The quotient is valid but generic: orbit size 384 and stabilizer size 1. BT1362 later upgrades this to a symmetric quotient with stabilizer \(C_2^4:C_4\) of order 64.}

\section{W(3,3) and the 240-Mode HoloNet Layer}

\claim{EXACT}{The symplectic polar graph \(\W\) has parameters \(\mathrm{SRG}(40,12,2,4)\). It has 40 vertices, 240 undirected edges, and 480 directed non-backtracking carriers.}

\claim{STRUCT}{The broader Q4 Diamond Machine paper interprets \(\W\) as the photonic substrate, with the Q4 router and toroidal heptad as complementary representations of the same finite architecture.}

\claim{EXACT}{The diamond identity from the broad synthesis is}
\[
40\cdot12\cdot3\cdot24\cdot15=518400=10\cdot 51840.
\]
This is arithmetic. Its physical interpretation remains structural until each bridge is separately certified.

\section{Toroidal Heptad, Reye, and Q4 Router}

\claim{STRUCT}{The parallel synthesis identifies a toroidal heptad of size 7, a Q4 router with 24 square plaquettes, and the Reye \((12_4,16_3)\) configuration as a shared geometric spine.}

\claim{EXACT}{The integer bridge}
\[
7\cdot24=168=6\cdot28
\]
is exact. The monodromy count
\[
18432=32\cdot24^2
\]
is also exact arithmetic.

\claim{SPEC}{The assertion that the heptad, Reye spine, and Q4 router are literally the same mathematical object in three photonic representations is promising but should remain structural/speculative until a functorial incidence-preserving map is machine-certified.}

\section{Kagome and Experimental Embedding}

\claim{STRUCT}{The Kagome route proposes that finite kagome or line-graph patches provide a physically accessible embedding of \(\W\). The key idea is to use flat-band localization and line-graph structure as a photonic realization path.}

\claim{SIM}{BT1342 turns the 40-mode/240-mode chip protocol into a synthetic falsifier simulator. It constructs \(\W\) over \(\Fthree^4\), verifies \(\mathrm{SRG}(40,12,2,4)\), and simulates pass/fail observables.}

\claim{SIM}{BT1345 replaces inherited synthetic phase targets with a matrix-derived standard Hashimoto calculation. The standard non-backtracking matrix gives phase clusters near}
\[
72.452^\circ \quad\text{and}\quad 127.087^\circ,
\]
not the earlier synthetic targets \(63.435^\circ\) and \(112.208^\circ\). Therefore the experimental protocol must either update the target angles or define a nonstandard normalized phase observable.

\section{Photonic Feasibility Gates}

\claim{ENG}{BT1335 shows that the 4320 base-channel layer is area-plausible under explicit chart-cell assumptions on a \(5\mathrm{mm}\times5\mathrm{mm}\) die.}

\claim{ENG}{BT1339 adds a first optical budget. In the conservative scenario, total loss is \(2.52\,\mathrm{dB}\), transmission is \(0.559\), and aggregate crosstalk is \(-24.21\,\mathrm{dB}\), passing the stated \(3\,\mathrm{dB}\) and \(-20\,\mathrm{dB}\) budgets.}

\claim{ENG}{This is not yet a foundry PDK simulation. Routed layout extraction, wavelength-dependent S-parameters, coupler fanout, thermal models, and packaging remain open.}

\section{Decoder and Threshold Boundary}

\claim{CERT}{BT1334 imposes the erasure-channel capacity wall}
\[
Q(p)=\max(0,1-2p),
\]
so a Gottesman--Knill decoder cannot push an erasure/loss threshold above \(50\%\).

\claim{SIM}{BT1336 gives the distance-only erasure benchmark for the \([[32,4,4]]\) block. It is a lower-bound/proxy benchmark, not a full ML or stabilizer decoder curve.}

\claim{CERT}{BT1341 and BT1362 now supply check objects for the next decoder benchmark: explicit \([[32,4,4]]\) gauge-quotient check pairs, including a symmetric clocked representative.}

\section{Physical Runtime Spine}

\claim{CERT}{BT1378 verifies the cross-artifact runtime contract: symmetric Q4 code, Q6/tomotope packet compiler, central \(C_3\) Steinberg scheduler, 2160-tick mirror bus, 8-tick packet word, and 51840-window Clifford runtime.}

% BT1380 post-BT1377 claim table
\begin{longtable}{rlll}
\toprule
BT & Class & Artifact & Role\\
\midrule
1362 & CERT & data/bt1362\_symmetric\_q4\_gauge\_quotient.json & Symmetric Q4 [[32,4,4]] quotient with $C_2^4:C_4$ clock\\
1363 & STRUCT & data/bt1363\_q4\_clock\_tomotope\_medial\_descent.json & Q4 clock descends to three tomotope sheets\\
1368 & CERT & data/bt1368\_q6\_tomotope\_equivariant\_flag\_lift.json & Q6/tomotope 192-flag equivariant lift\\
1374 & CERT & data/bt1374\_q6\_tomotope\_packet\_route\_compiler.json & Packet route compiler to Q6 edge addresses\\
1375 & CERT & data/bt1375\_steinberg\_cycle\_operator\_scheduler\_lift.json & Central $C_3$ Steinberg scheduler\\
1376 & CERT & data/bt1376\_s3\_gauge\_radius3\_local\_optimum\_certificate.json & Radius-3 local optimum for S3 gauge\\
1377 & ENG & data/bt1377\_physical\_universal\_computation\_contract.json & Physical Clifford runtime plus non-Clifford port boundary\\
1378 & CERT & data/bt1378\_runtime\_contract\_verification.json & Cross-artifact runtime contract verifier\\
1379 & OPT & data/bt1379\_s3\_gauge\_max2csp\_spec.json & Global S3 gauge Max-2CSP frontier\\
\bottomrule
\end{longtable}


% BT1381--BT1383 runtime frontier insert
\section{Post-BT1377 Runtime Frontier}

The BT1362--BT1380 stack gives a certified physical Clifford/symplectic runtime.  BT1381--BT1383 extend that stack in three directions: global S3 gauge optimization, explicit non-Clifford port ABI, and paper/release integration.

\subsection{S3 gauge optimization frontier}

BT1379 frames synchronization as a Max-2CSP on 40 W33 line variables, six S3 labels per variable, and 540 skew-line constraints.  BT1381 adds a deterministic global solver probe.  With seed 1381 and 240 coordinate-ascent restarts, no witness exceeds the BT1373 score of 210 identity residual edges.  The frontier remains:

\[
\mathrm{identity}=210,\qquad \mathrm{corrections}=330.
\]

This is not a global optimality proof.  It is a solver-frontier certificate: branch-and-bound, ILP/SAT, or SDP relaxation remains the next proof target.

\subsection{Non-Clifford port ABI}

BT1382 certifies the non-Clifford port as an architecture boundary.  The deterministic runtime is a protected Clifford machine of order 51840 and is not universal by itself.  Universal computation requires one of the explicit ABI ports:

\begin{itemize}
\item Hesse-SIC/T measurement port;
\item Fibonacci braiding port.
\end{itemize}

The port contract records entry boundary, packet interface, and verification signature.  It does not certify a magic-state factory, Hesse-SIC hardware stack, or Fibonacci anyon device.

\subsection{Claim-class integration}

The runtime frontier is claim-stratified as follows:

\begin{center}
\begin{tabular}{rll}
\toprule
BT & Class & Role\\
\midrule
1378 & CERT & cross-artifact runtime contract verifier\\
1379 & OPT & S3 gauge Max-2CSP specification\\
1381 & SIM/OPT & deterministic global solver probe\\
1382 & ENG/ABI & non-Clifford port ABI certificate\\
1383 & CERT & paper and release integration manifest\\
\bottomrule
\end{tabular}
\end{center}


\section{Hesse-SIC/T Non-Clifford Port}

\claim{ABI}{BT1382 records two non-Clifford port options. BT1385 makes the Hesse-SIC/T option concrete as a qutrit SIC token with nine measurement outcomes.}

The port token is consumed at a Clifford packet boundary. The measurement signature records an outcome
\[
h\in\{0,1,\ldots,8\},
\]
maps it to a Clifford correction plus one T-frame parity bit, appends a correction side record, and returns to the Clifford packet ABI. Failure modes are missing outcome, ambiguous SIC outcome, feed-forward timeout, and Clifford ABI restoration failure.

\claim{ENG}{This is an ABI certificate, not a physical resource-factory certificate. Hesse-SIC optics, magic-state yield, and threshold overhead remain future gates.}

\section{Speculative Bridges Held Apart}

\claim{SPEC}{The Monster moonshine bridge around \(14641=11^4\), the Calabi--Yau threefold branch, and the W99 syndrome-holonet cousin are valuable idea generators. They should be treated as speculative or structural until executable finite certificates are built.}

\claim{STRUCT}{The strongest of these is the W99 clue: \(1620=540\cdot3\), matching the independent syndrome count. This suggests a self-similar syndrome layer and deserves a dedicated finite audit.}

\section{Conclusion}

The merged Q4 Diamond Machine is no longer a single undifferentiated claim. Its exact core is audited. Its \([[32,4,4]]\) code now has concrete gauge-quotient certificates, including a symmetric clocked representative. Its W(3,3) graph layer has both a synthetic falsifier simulator and a matrix-derived Hashimoto correction. Its runtime layer has a verified Clifford packet contract, a precise S3 optimization frontier, and an explicit non-Clifford ABI boundary. Its photonic layout claims remain engineering gates, and its Monster, Calabi--Yau, W99, and deeper unification bridges remain clearly marked as future certificate targets.

\end{document}
